\chapter[Thảo luận, kết luận và hướng phát triển]{Thảo luận, kết luận\\và hướng phát triển}
\label{chap:discussion}

\section{Thảo luận các phát hiện chính}
Chuỗi thực nghiệm ở Chương~\ref{chap:results} làm sáng tỏ ba bài học thiết kế đáng chú ý cho bài toán học sâu sống còn đa phương thức trên dữ liệu y tế cỡ mẫu vừa và nhỏ.

Thứ nhất, \emph{hiệu quả nén tham số của phân rã CP cần được nhìn đúng theo từng phạm vi kiến trúc}. Thay lớp chiếu tích ngoài bằng phân rã CP giúp giảm tới \textbf{98.38\%} tham số ở ba lớp song tuyến đầu (từ 229696 xuống 3712) và 90.90\% khi tính cả khối MLP tương tác, nhưng chỉ giảm khoảng 0.67\% tổng tham số toàn mạng vì bộ mã hóa ảnh ResNet-18 3D chiếm phần áp đảo. Điều này có nghĩa là lợi ích lớn nhất của CP trong kiến trúc này không nằm ở việc thu nhỏ dung lượng toàn mạng, mà ở việc \emph{tránh quá khớp cục bộ tại tầng hợp nhất} — nơi dễ phình trọng số nhất khi kết hợp nhiều nguồn đặc trưng trên 251 bệnh nhân huấn luyện. Mức đánh đổi C-index giữa Full CP và Full Outer-Product chỉ là 0.0066, là cái giá hợp lý cho sự tinh gọn ở tầng tương tác.

Thứ hai, \emph{nhiều cặp tương tác hơn không đồng nghĩa với kết quả tốt hơn}. Trong bảy cấu trúc CP có tương tác cặp, RC-Free chỉ giữ hai cặp MR và MC với bốn nhánh phần dư nhưng đạt C-index xác thực trung bình cao nhất (\textbf{0.8347}). Kết quả này nhỉnh hơn Full CP (0.8278) và giảm được 15460 tham số; thậm chí cấu trúc một cặp RC-only cũng đạt 0.8319, cao hơn Full CP. Vì cả ba nguồn dữ liệu vẫn giữ đường truyền gốc qua phần dư có cổng, việc cộng thêm cả ba cặp tương tác có thể gây dư thừa và tăng phương sai ước lượng trên mẫu nhỏ. Khi đặt ảnh MRI làm cầu nối sang radiomics và lâm sàng, RC-Free là cấu trúc gọn nhẹ và hiệu quả nhất.

Thứ ba, \emph{ý nghĩa thực tiễn của phân tầng nguy cơ luôn gắn với cỡ mẫu và số biến cố}. Phân tích Kaplan--Meier cho thấy điểm nguy cơ của RC-Free tách biệt hai nhóm diễn biến: trên tập xác thực đạt log-rank $p=0.0134$ (12/26 so với 4/28 biến cố) và trên tập kiểm tra hậu nghiệm đạt $p<0.001$. Tuy nhiên, khi xét theo từng thời hạn, AUC xác thực tại 2 và 3 năm (0.9805 và 0.9788) chỉ dựa trên 7 biến cố tích lũy. Tại mốc 5 năm, AUC giảm xuống 0.7889 kèm độ lệch chuẩn rộng hơn khi số người còn theo dõi quá 5 năm chỉ còn 19 (Bảng~\ref{tab:horizon-support}).

\section{Giới hạn nghiên cứu}
\label{sec:study-limitations}
Nghiên cứu còn tồn tại bốn giới hạn chính sau:
\begin{enumerate}[itemsep=0.15em,topsep=0.25em]
\item \textbf{Dữ liệu mới dừng ở bộ ADNI, chưa thử nghiệm tại Việt Nam.} Toàn bộ quá trình huấn luyện và đánh giá hiện chỉ sử dụng bộ dữ liệu quốc tế ADNI với tập kiểm chứng còn khiêm tốn (54 người xác thực và 54 người kiểm tra hậu nghiệm). Mô hình chưa được thu thập và kiểm thử trên dữ liệu thực tế từ các bệnh viện lớn ở Việt Nam, nơi thiết bị chụp MRI, quy trình đo lâm sàng và đặc điểm người bệnh có thể khác biệt.
\item \textbf{Thiết kế mốc thời gian hồi cứu.} Thời gian sống còn $t_i$ được tính từ ngày chụp MRI thứ ba $p_i$, nhưng quy tắc lọc mẫu lại giữ những người chưa tiến triển đến hết mốc 18 tháng $L_i$ (nên có nguy cơ sai lệch thời gian bất tử nếu diễn giải như dự báo ngay tại ngày chụp $p_i$ \cite{suissa2008immortal}). Ngoài ra, cửa sổ ghép $\pm90$ ngày làm 57 người có điểm MMSE hoặc CDR gắn với lần chụp cuối thực tế được đo sau ngày $p_i$ (Mục~\ref{sec:data-split}), khiến kết quả hiện tại phù hợp hơn cho đánh giá trên nhóm hồi cứu.
\item \textbf{Thang đo đầu vào và định nghĩa biến cố cùng dùng MMSE và CDR.} Dù hai giai đoạn đo được tách biệt trước và sau mốc 18 tháng, việc cùng dựa vào ngưỡng điểm MMSE và CDR để xác định biến cố vẫn tạo sự phụ thuộc vào hai thang đo này; mô hình chưa được đối chiếu với chẩn đoán lâm sàng độc lập hay các xét nghiệm sinh học chuyên sâu hơn.
\item \textbf{Đầu ra mới dừng ở mức xếp hạng nguy cơ.} Điểm $\eta_i$ của mô hình Cox cho biết ai có nguy cơ tiến triển sớm hơn ai, nhưng chưa quy đổi trực tiếp thành xác suất tiến triển cụ thể theo phần trăm cho từng bệnh nhân và chưa đo tốc độ xử lý thực tế trên máy tính tại bệnh viện.
\end{enumerate}

\section{Hướng phát triển}
Từ các giới hạn trên, khóa luận có thể mở rộng theo ba hướng thực tế:
\begin{enumerate}[itemsep=0.15em,topsep=0.25em]
\item \emph{Thử nghiệm trên dữ liệu bệnh viện tại Việt Nam}: Phối hợp thu thập dữ liệu theo dõi dọc từ các bệnh viện lớn trong nước để kiểm chứng, hiệu chỉnh mô hình phù hợp với bệnh nhân Việt Nam và bổ sung thêm các nguồn thông tin xét nghiệm sẵn có.
\item \emph{Hoàn thiện thiết kế dự báo theo mốc thời gian}: Giới hạn chặt dữ liệu lâm sàng đến trước ngày chụp cuối $p_i$ hoặc áp dụng thiết kế dự báo tại mốc cố định $L_i$ \cite{vanhouwelingen2007dynamic} để loại bỏ hoàn toàn độ lệch thời gian giữa đầu vào và kết cục.
\item \emph{Quy đổi xác suất cá thể hóa và hỗ trợ giải thích}: Ước lượng thêm hàm sống còn nền \cite{kvamme2019time} để chuyển điểm nguy cơ thành xác suất tiến triển theo từng năm, kết hợp phân tích trọng số attention và cổng điều tiết nhằm giúp bác sĩ dễ theo dõi cơ sở dự báo và đo lường thời gian chạy thực tế khi triển khai ứng dụng.
\end{enumerate}

\section{Kết luận}
Tiên lượng tiến triển ở người mắc suy giảm nhận thức nhẹ đòi hỏi không chỉ theo dõi diễn biến qua nhiều lần khám, mà còn kết hợp hài hòa giữa hình ảnh cấu trúc não, đặc trưng định lượng và đánh giá lâm sàng. Trên tập dữ liệu hồi cứu gồm 359 người tham gia ADNI với ba mốc thăm khám, khóa luận đã xây dựng một quy trình phân tích sống còn đa phương thức hoàn chỉnh tích hợp ảnh MRI 3D, radiomics toàn não và thang điểm MMSE/CDR. Các thực nghiệm đối chứng khẳng định từng bước đóng góp của thiết kế: chuỗi MRI dọc với T-LSTM và attention vượt trội so với ảnh đơn thời điểm, kết hợp đủ ba nguồn nâng cao rõ rệt năng lực xếp hạng nguy cơ, và mô hình hóa tương tác song tuyến theo cặp khai thác mối liên hệ giữa các nguồn hiệu quả hơn phép ghép nối thông thường.

Kiến trúc \textbf{RC-Free} kết hợp hai bước tinh giản: phân rã tensor CP hạng thấp giúp giảm 98.38\% tham số ở ba lớp song tuyến đầu, và lược bỏ cặp tương tác trực tiếp RC giúp tiết kiệm thêm 15460 tham số so với Full CP mà vẫn giữ đủ ba nguồn đầu vào. Trên tập xác thực, RC-Free đạt C-index trung bình cao nhất trong bảy cấu trúc CP ($0.8347\pm0.0237$) và phân tầng có ý nghĩa thống kê hai nhóm nguy cơ sống còn ($p=0.0134$). Kết quả này cho thấy kết hợp tham số hóa hạng thấp với chọn lọc cặp tương tác hợp lý là hướng đi gọn nhẹ và khả thi cho tiên lượng đa phương thức trên dữ liệu y tế cỡ mẫu hạn chế.
